\section{Reusable algebraic and functional-analytic infrastructure}

\begin{lemma}[Quotient by the kernel]\label{lem:quotient-range}
Let $X$ and $Y$ be real Banach spaces and let $A:X\to Y$ be bounded and linear. The map
induced by $A$ is a linear isomorphism
\[
  X/\ker A\simeq\ran A,\qquad [x]\longmapsto Ax.
\]
Consequently, if $A$ has finite rank, then $\ker A$ is closed, $X/\ker A$ is finite-dimensional,
and
\[
  \codim_X\ker A=\dim\ran A=\rank A.
\]
\uses{def:fredholm-rank}
\lean{KaltonPeck.Support.Fredholm.quotientKernelRangeEquiv, KaltonPeck.Support.Fredholm.quotientKernelRangeEquiv_apply_mk, KaltonPeck.Support.Fredholm.quotientKernelRange}
\end{lemma}

\begin{proof}
Send the class of $x$ to $Ax$. This is well-defined because two representatives differ by an
element of the kernel. It is injective by the same observation and surjective by the definition
of the range. Continuity of $A$ makes its kernel closed. Transporting finite-dimensionality and
dimension through the isomorphism gives the last assertion.
\end{proof}

\begin{lemma}[Continuous double annihilator]\label{lem:double-annihilator}
Let $W$ be a closed subspace of a real Banach space $X$ and let
\[
  \Ann(W)=\{f\in\strongdual X:f(w)=0\text{ for every }w\in W\}.
\]
Then
\[
  W=\{x\in X:f(x)=0\text{ for every }f\in\Ann(W)\}.
\]
\lean{KaltonPeck.Support.Forms.continuousAnnihilator, KaltonPeck.Support.Forms.continuousDoubleAnnihilator}
\end{lemma}

\begin{proof}
Every $w\in W$ is annihilated by definition. Conversely, if $x\notin W$, the continuous
Hahn--Banach separation theorem supplies $f\in\strongdual X$ and a strict separation between $x$
and $W$. The functional must vanish on $W$: if $f(w)\ne0$ for some $w\in W$, arbitrary real
multiples of $w$ make $f$ unbounded above and below on $W$, contradicting strict separation.
The separating inequality then gives $f(x)\ne0$, so $x$ does not belong to the displayed double
annihilator.
\end{proof}

\begin{lemma}[Continuous dual of a quotient]\label{lem:quotient-dual}
If $W$ is closed in a Banach space $X$, pullback along the quotient map $q:X\to X/W$ is a
continuous linear isomorphism
\[
  \strongdual{(X/W)}\simeq\Ann(W),
  \qquad \varphi\longmapsto\varphi\circ q.
\]
\lean{KaltonPeck.Support.Forms.quotientDualEquivAnnihilator, KaltonPeck.Support.Forms.quotientDualEquivAnnihilator_apply}
\end{lemma}

\begin{proof}
Pullback is injective because $q$ is onto and its image lies in $\Ann(W)$. If
$f\in\Ann(W)$, the rule $[x]\mapsto f(x)$ is well-defined and bounded for the quotient norm, and
is inverse to pullback. Since $W$ is closed, $X/W$ is Banach, while $\Ann(W)$ is a closed
subspace of $\strongdual X$ and hence Banach. The bounded inverse theorem therefore promotes
pullback to a continuous linear isomorphism.
\end{proof}

\begin{lemma}[Strong symplectic spaces are reflexive]\label{lem:strong-reflexive}
Every strong symplectic Banach space is reflexive.
\uses{def:strong-form,def:transpose-adjoint}
\lean{KaltonPeck.Support.Forms.strongSymplecticTransposeInclusion, KaltonPeck.Support.Forms.strongSymplecticReflexive}
\end{lemma}

\begin{proof}
Let $D:X\to\strongdual X$ be the isomorphism induced by the form and let
$\iota:X\to\strongdual{\strongdual X}$ be the canonical embedding. Alternation gives
\[
  D^*\iota=-D.
\]
Indeed, $(D^*\iota x)(y)=\omega(y,x)=-\omega(x,y)$. The transpose
$D^*:\strongdual{\strongdual X}\to\strongdual X$ is a continuous linear isomorphism. Hence
$\iota=-(D^*)^{-1}D$ is itself a continuous linear isomorphism, in particular onto, which is
exactly reflexivity.
\end{proof}

\begin{lemma}[Symplectic-adjoint calculus]\label{lem:adjoint-calculus}
For bounded operators $A,B$ on a strong symplectic Banach space and $a\in\R$,
\[
 I^+=I,\qquad (A+B)^+=A^++B^+,\qquad (aA)^+=aA^+,
\]
\[
 (AB)^+=B^+A^+,\qquad (A^+)^+=A,\qquad
 \omega(A^+x,y)=\omega(x,Ay).
\]
\uses{def:transpose-adjoint,lem:strong-reflexive}
\lean{KaltonPeck.Support.Forms.adjoint_one, KaltonPeck.Support.Forms.adjoint_add, KaltonPeck.Support.Forms.adjoint_smul, KaltonPeck.Support.Forms.adjoint_mul, KaltonPeck.Support.Forms.adjoint_involutive, KaltonPeck.Support.Forms.adjoint_apply}
\end{lemma}

\begin{proof}
The defining equation gives the last identity. Identity, addition, and scalar compatibility follow
from the corresponding transpose laws. Transposition reverses composition, giving the product
law. Apply the defining equation twice; alternation and nondegeneracy of $\omega$ then imply
$(A^+)^+=A$. Equivalently, one may expand $D^{-1}A^*D$ and use the canonical double-dual map,
whose surjectivity is Lemma~\ref{lem:strong-reflexive}.
\end{proof}

\begin{lemma}[Fredholm calculus used below]\label{lem:fredholm-calculus}
Let $X,Y,Z,X',Y'$ be real Banach spaces.
\begin{enumerate}
\item If $A:X\to Y$ is bounded and $U:Y\simeq Y'$ and $V:X'\simeq X$ are bounded linear
isomorphisms, then $A$ is Fredholm if and only if $UAV$ is Fredholm. The induced kernel and
cokernel isomorphisms preserve nullity, cokernel dimension, and Fredholm index.
\item Multiplication by a nonzero real scalar preserves Fredholmness and index. If
$A:X\to Y$ and $C:Y\to Z$ are bounded and $C$ is injective, then
\[
  \ker(CA)=\ker A.
\]
\item If $A:X\to Y$ and $B:Y\to Z$ are Fredholm, then $BA$ is Fredholm and
\[
  \operatorname{ind}(BA)=\operatorname{ind}B+\operatorname{ind}A.
\]
\item If a bounded map $A:X\to Y$ has a bounded left inverse, then $\ker A=0$ and $\ran A$ is
closed; hence it has finite-dimensional kernel and closed range.
\end{enumerate}
\uses{def:fredholm-rank}
\lean{KaltonPeck.Support.Fredholm.isFredholm_equiv_comp, KaltonPeck.Support.Fredholm.kernelEquivOfEquivComp, KaltonPeck.Support.Fredholm.cokernelEquivOfEquivComp, KaltonPeck.Support.Fredholm.nullity_equiv_comp, KaltonPeck.Support.Fredholm.cokernelFinrank_equiv_comp, KaltonPeck.Support.Fredholm.fredholmIndex_equiv_comp, KaltonPeck.Support.Fredholm.fredholm_smul, KaltonPeck.Support.Fredholm.ker_comp_of_injective, KaltonPeck.Support.Fredholm.isFredholm_comp, KaltonPeck.Support.Fredholm.fredholmIndex_comp, KaltonPeck.Support.Fredholm.leftInverseKernelRange}
\end{lemma}

\begin{proof}
Bounded linear isomorphisms carry kernels onto kernels, ranges onto ranges, and the corresponding
quotients onto one another; they also preserve closedness and dimension. This proves the first
assertion. Multiplication by a nonzero scalar leaves both kernel and range unchanged. If $C$ is
injective, then $CAx=0$ if and only if $Ax=0$, proving the displayed kernel identity.

For composition, $\ran A+\ker B$ is closed because $\ran A$ is closed and $\ker B$ is
finite-dimensional. Its image in $Y/\ker B$ is closed. The map induced by $B$ is a bounded
linear isomorphism
\[
  Y/\ker B\simeq\ran B,
\]
so it carries that image onto $B(\ran A)=\ran(BA)$. Since $\ran B$ is closed in $Z$,
$\ran(BA)$ is closed in $Z$.

There is an exact sequence
\[
0\longrightarrow\ker A\longrightarrow\ker(BA)\longrightarrow\ker B
\longrightarrow Y/\ran A\longrightarrow Z/\ran(BA)\longrightarrow Z/\ran B
\longrightarrow0.
\]
The map $\ker(BA)\to\ker B$ is induced by $A$; the next map sends $y\in\ker B$ to
$y+\ran A$; the map $Y/\ran A\to Z/\ran(BA)$ is induced by $B$; and the last map is the
natural quotient map. Their kernels are respectively the images of the preceding maps:
$\ker A$, $\ker B\cap\ran A$, $(\ran A+\ker B)/\ran A$, and
$\ran B/\ran(BA)$. The last map is surjective. Thus every term is finite-dimensional, and
the alternating dimension identity for this exact sequence gives Fredholmness of $BA$ and
index additivity.

Finally, if $LA=I$, then $A$ is injective and
\[
  \|x\|=\|LAx\|\leq\|L\|\,\|Ax\|.
\]
Consequently a convergent sequence in $\ran A$ has Cauchy, hence convergent, preimages; its limit
maps to the given range limit. Thus $\ran A$ is closed.
\end{proof}

\begin{lemma}[Range of a skew Fredholm form]\label{lem:skew-fredholm-range}
Let $X$ be reflexive and let $S:X\to\strongdual X$ be Fredholm. Suppose that
$S^*\iota_X=-S$, where $\iota_X:X\to\strongdual{\strongdual X}$ is the canonical embedding. If
$N=\ker S$, then $S$ has Fredholm index zero and
\[
  \ran S=\Ann(N).
\]
\uses{def:transpose-adjoint,def:fredholm-rank,lem:double-annihilator,lem:quotient-dual}
\lean{KaltonPeck.Support.Fredholm.skewFredholmRange}
\end{lemma}

\begin{proof}
Since $S$ is Fredholm, $N$ is finite-dimensional and $\ran S$ is closed. The skew identity and
surjectivity of $\iota_X$ give
\[
  \ker S^*=\iota_X(N).
\]
This kernel is exactly the continuous annihilator of $\ran S$ in $\strongdual{\strongdual X}$.
Applying Lemma~\ref{lem:double-annihilator} to the closed subspace
$\ran S\subseteq\strongdual X$ therefore gives
\[
  \ran S=\{f\in\strongdual X:f(n)=0\text{ for every }n\in N\}=\Ann(N).
\]
Lemma~\ref{lem:quotient-dual}, applied to $\ran S\subseteq\strongdual X$, identifies the dual of
$(\strongdual X)/\ran S$ with $\ker S^*=\iota_X(N)$. The cokernel is finite-dimensional, so it and
its dual have the same dimension. Thus the cokernel dimension equals $\dim N$, and the Fredholm
index is zero.
\end{proof}

\begin{lemma}[Even rank of a finite alternating form]\label{lem:finite-alt-even}
Let $V$ be a finite-dimensional real vector space and let $c$ be an alternating bilinear form on
$V$. Put $\rad c=\{x\in V:c(x,y)=0\text{ for every }y\in V\}$. The induced map
$V\to V^*$ has even rank. Consequently,
\[
  \dim\rad c\equiv\dim V\pmod2.
\]
\lean{KaltonPeck.Support.FiniteParity.finiteAlternatingRankEven, KaltonPeck.Support.FiniteParity.finiteContinuousAlternatingRankEven}
\end{lemma}

\begin{proof}
Induct on $\dim V$. If $c=0$, its rank is zero. Otherwise choose $u,v$ with $c(u,v)\ne0$.
Put $a=c(u,v)$ and $P=\operatorname{span}\{u,v\}$. For the fixed form $c$, write
\[
  P^{\circ_c}=\{z\in V:c(z,p)=0\text{ for every }p\in P\},
\]
and abbreviate this subspace by $P^\circ$ inside this proof. The restriction to $P$ is
nondegenerate. For $x\in V$, define
\[
  p(x)=\frac{c(x,v)}a\,u-\frac{c(x,u)}a\,v.
\]
Then $x-p(x)\in P^\circ$, so $V=P\oplus P^\circ$. Both cross terms vanish, and the associated
map is the direct sum of a rank-two map on $P$ and the map of the alternating restriction to
$P^\circ$. Since $\dim P^\circ=\dim V-2$, induction makes the latter rank even. Thus the total
rank is even, and rank--nullity gives the congruence.
\end{proof}

\begin{lemma}[Finite complex structures have even dimension]\label{lem:finite-complex-even}
If a finite-dimensional real vector space $V$ has an endomorphism $J$ with $J^2=-I$, then
$\dim_{\R} V$ is even.
\lean{KaltonPeck.Support.FiniteParity.finiteComplexDimensionEven, KaltonPeck.Support.FiniteParity.finiteContinuousComplexDimensionEven, KaltonPeck.Support.FiniteParity.finiteInvariantSubmoduleComplexDimensionEven}
\end{lemma}

\begin{proof}
Define multiplication by $a+ib\in\mathbb C$ on $V$ by
\[
  (a+ib)v=av+bJv.
\]
The equation $J^2=-I$ verifies the complex scalar axioms, and the restricted real scalar action is
the original one. A finite real spanning set also spans over $\mathbb C$, so $V$ is
finite-dimensional over $\mathbb C$. The scalar-tower dimension formula therefore gives
\[
  \dim_{\R} V=\dim_{\R}\mathbb C\,\dim_{\mathbb C}V=2\dim_{\mathbb C}V.
\]
\end{proof}
